\section{Tujuan Buku Ajar}

Buku ajar ini dirancang sebagai panduan komprehensif untuk menguasai \crypt{Kriptografi} dalam konteks Programme Studi Matematika S1 \cite{stinson2018}. Fokus utama buku ini adalah pada pemahaman matematis, konsep teoretis, dan implementasi algoritma kriptografi menggunakan bahasa pemrograman C. Tujuan spesifik buku ini adalah:

\begin{enumerate}
  \item Memberikan pemahaman mendalam tentang prinsip-prinsip dasar kriptografi dan keamanan informasi
  \item Mengembangkan kemampuan menerapkan matematika (aritmetika modular, teori bilangan) dalam konteks kriptografi
  \item Membangun keterampilan implementasi algoritma kriptografi klasik dan modern dalam bahasa C
  \item Menumbuhkan kemampuan analisis kriptanalisis terhadap cipher sederhana
  \item Memfasilitasi pencapaian CPL dan CPMK yang telah ditetapkan dalam RPS berbasis OBE
\end{enumerate}

Setelah mempelajari buku ini secara menyeluruh, mahasiswa diharapkan mampu:
\begin{itemize}
  \item Menjelaskan konsep dasar kriptografi (confidentiality, integrity, authenticity)
  \item Menerapkan aritmetika modular dan algoritma Euclid dalam konteks kriptografi
  \item Mengimplementasikan cipher Caesar, Vigenère, dan RSA dalam bahasa C
  \item Menganalisis kelemahan cipher klasik melalui frequency analysis dan Kasiski examination
  \item Merancang dan menganalisis protokol kriptografi dasar
\end{itemize}

Pendidikan kriptografi modern membutuhkan fondasi matematis yang kuat termasuk teori bilangan, aljabar linear, dan matematika diskrit untuk memahami algoritma enkripsi yang kompleks. Menurut Coursera, pemahaman tentang modular arithmetic, algoritma Euclid, dan teorema sisa China merupakan prasyarat esensial untuk menguasai kriptografi tingkat lanjut. Kurikulum yang dirancang dengan baik mengintegrasikan konsep matematika murni dengan aplikasi praktis dalam keamanan digital, mempersiapkan mahasiswa untuk karir di bidang keamanan siber.

Profesi kriptografer menuntut kombinasi unik antara keahlian matematika, pemrograman, dan pemahaman keamanan siber untuk mengembangkan cipher baru dan melindungi data sensitif. Industri saat ini sangat membutuhkan ahli kriptografi yang mampu membuat solusi keamanan inovatif seiring dengan evolusi ancaman digital yang terus berubah. Mahasiswa yang menguasai materi dalam buku ini akan memiliki kompetensi untuk berkarir di lembaga pemerintah, institusi keuangan, atau perusahaan teknologi yang memprioritaskan keamanan data.
